%! Inverse Functions — Unit 6, Lesson 6.7 — Homework (Answer Key)
\documentclass[10pt]{article}
\usepackage{algebra2-article}
\usepackage{algebra2-key}   % pulls in -boxes; do NOT also load -boxes
\newcommand{\TallMath}[1]{$\displaystyle #1\rule[-1.4em]{0pt}{3.2em}$}
% In-formula answer slot (\ans is text-mode and must never appear inside $...$).
\newcommand{\mans}[1]{{\color{keyred}\mathbf{#1}}}
\newcommand{\mexp}[1]{{\color{keyred}#1}}
\newcommand{\lgrid}[4]{%
  \draw[step=1, linegray!40, very thin] (#1,#3) grid (#2,#4);
  \draw[->, semithick, charcoal] (#1,0)--(#2,0) node[below right, font=\scriptsize]{$x$};
  \draw[->, semithick, charcoal] (0,#3)--(0,#4) node[left, font=\scriptsize]{$y$};}

\begin{document}

\pageheader{Unit 6, Lesson 6.7}{Homework --- Answer Key}
\namedateperiod

\vspace{0.10in}

\begin{notesbox}{Practice}
Reverse the operations \emph{and} the order. Then test what you built --- both ways.

\begin{enumerate}[leftmargin=1.9em, itemsep=11pt]
  \item \textbf{Are they inverses?} Compose both ways before you answer. One of these four pairs is
        an impostor.
    \begin{center}
    \renewcommand{\arraystretch}{1.9}
    \begin{tabularx}{\linewidth}{@{}l X X c@{}}
    \hline
    \textbf{$f$ and $g$} & \textbf{$f\big(g(x)\big)$} & \textbf{$g\big(f(x)\big)$} & \textbf{Inverses?} \\
    \hline
    (a) $f(x)=x-9$, \; $g(x)=x+9$ & \ans{$(x+9)-9=x$} & \ans{$(x-9)+9=x$} & \ans{yes} \\
    (b) $f(x)=\dfrac{x}{2}$, \; $g(x)=2x$ & \ans{$\frac{2x}{2}=x$} & \ans{$2\!\left(\frac{x}{2}\right)=x$} & \ans{yes} \\
    (c) $f(x)=5x-2$, \; $g(x)=\dfrac{x}{5}+2$ & \ans{$5\!\left(\frac{x}{5}+2\right)-2=x+8$} & \ans{$\frac{5x-2}{5}+2=x+\frac{8}{5}$} & \ans{no} \\
    (d) $f(x)=x^{3}+2$, \; $g(x)=\sqrt[3]{x-2}$ & \ans{$\left(\sqrt[3]{x-2}\right)^{3}+2=x$} & \ans{$\sqrt[3]{(x^{3}+2)-2}=x$} & \ans{yes} \\
    \hline
    \end{tabularx}
    \end{center}
    \vspace{0.10in}
    (e) Repair the impostor: the true inverse is \ans{$\frac{x+2}{5}$} \\[2pt]
    \hspace*{0.3in}\emph{What exactly did that row get wrong --- the operations, or the order?}
    \ans{the order (add $2$ \emph{first}, then divide)}

  \item \textbf{Swap and solve.} Two of these six need a restriction. Write it if it is needed.
    \begin{center}
    \renewcommand{\arraystretch}{1.9}
    \begin{tabularx}{\linewidth}{@{}l X l X@{}}
    \hline
    \textbf{$f(x)$} & \textbf{$f^{-1}(x)$} & \textbf{$f(x)$} & \textbf{$f^{-1}(x)$} \\
    \hline
    (a) $x-11$ & \ans{$x+11$} & (d) $x^{3}-8$ & \ans{$\sqrt[3]{x+8}$} \\
    (b) $\dfrac{x}{3}+6$ & \ans{$3x-18$} & (e) $\sqrt[3]{x+1}$ & \ans{$x^{3}-1$} \\
    (c) $7-2x$ & \ans{$\frac{7-x}{2}$} & (f) $\sqrt{x}+4$ & \ans{$(x-4)^{2}$, $x\ge4$} \\
    \hline
    \end{tabularx}
    \end{center}
    \vspace{2pt}
    {\small \emph{Work, the two that bite:} \; (c) $x=7-2y \Rightarrow 2y=7-x \Rightarrow
    y=\frac{7-x}{2}$ --- equivalently $-\frac12x+\frac72$; this one is its own kind of trap, since
    the $-2$ has to be undone \emph{after} the $7$ is moved. \;
    (f) $x=\sqrt{y}+4 \Rightarrow x-4=\sqrt{y} \Rightarrow y=(x-4)^{2}$, and $\sqrt{y}\ge0$ forces
    $x-4\ge0$, i.e.\ $x\ge4$. \; Row (e) needs nothing: $x=\sqrt[3]{y+1}\Rightarrow x^{3}=y+1$.}
    \vspace{4pt}

    Rows (d) and (e) both involve cubes, and neither needed a restriction. Row (f) did. What is it
    about an \emph{even} index that makes the difference? \ans{An even root never returns a negative
    number, so $f$'s range is only half a line --- and that half is all $f^{-1}$ may accept.}

  \item \textbf{From a table.} Here are some values of $g$.
    \begin{center}
    \renewcommand{\arraystretch}{1.4}
    \begin{tabular}{@{}|c|c|c|c|c|c|@{}}
    \hline
    $x$ & $-2$ & $0$ & $1$ & $3$ & $4$ \\
    \hline
    $g(x)$ & $-3$ & $1$ & $3$ & $7$ & $9$ \\
    \hline
    \end{tabular}
    \hspace{0.25in}
    \renewcommand{\arraystretch}{1.4}
    \begin{tabular}{@{}|c|c|c|c|c|c|@{}}
    \hline
    $x$ & \ans{$-3$} & \ans{$1$} & \ans{$3$} & \ans{$7$} & \ans{$9$} \\
    \hline
    $g^{-1}(x)$ & \ans{$-2$} & \ans{$0$} & \ans{$1$} & \ans{$3$} & \ans{$4$} \\
    \hline
    \end{tabular}
    \end{center}
    \vspace{-2pt}
    (a) $g^{-1}(7)=$ \ans{$3$} \hspace{0.25in}
    (b) $g^{-1}(1)=$ \ans{$0$} \hspace{0.25in}
    (c) $g\big(g^{-1}(9)\big)=$ \ans{$9$} \\[3pt]
    (d) The graph of $g$ also passes through $(5,11)$. Name a point that must be on the graph of
    $g^{-1}$: \ans{$(11,5)$} \\[3pt]
    (e) These values follow the rule $g(x)=2x+1$. Find $g^{-1}(x)$ and check it against your table.
    \; $g^{-1}(x)=$ \ans{$\frac{x-1}{2}$}

  \item \textbf{From two graphs.} Below are $y=h(x)$ and $y=k(x)$. The dashed line is $y=x$.
    \begin{center}
    \begin{minipage}[c]{0.47\linewidth}
    \centering
    \begin{tikzpicture}[scale=0.34]
      \lgrid{-5}{6}{-5}{5}
      \draw[dashed, semithick, charcoal] (-4.6,-4.6)--(4.6,4.6);
      \draw[very thick, forest] (-4,-3)--(-1,0)--(5,3);
      \fill[forest] (-4,-3) circle (3.4pt);
      \fill[forest] (-1,0) circle (3.4pt);
      \fill[forest] (1,1) circle (3.4pt);
      \fill[forest] (5,3) circle (3.4pt);
      \node[forest, font=\scriptsize] at (-3.0,2.4) {$y=h(x)$};
    \end{tikzpicture}
    \end{minipage}
    \hfill
    \begin{minipage}[c]{0.47\linewidth}
    \centering
    \begin{tikzpicture}[scale=0.34]
      \lgrid{-5}{5}{-5}{5}
      \begin{scope}
        \clip (-5,-5) rectangle (5,5);
        \draw[very thick, navy, domain=-2.9:2.9, samples=120, smooth] plot (\x,{(\x)*(\x)-3});
      \end{scope}
      \node[navy, font=\scriptsize] at (-3.2,-4.2) {$y=k(x)$};
    \end{tikzpicture}
    \end{minipage}
    \end{center}
    \vspace{-2pt}
    (a) Four points are marked on $h$. List the four matching points on $h^{-1}$: \\[3pt]
    \hspace*{0.3in} \ans{$(-3,-4)$} \hspace{0.15in} \ans{$(0,-1)$} \hspace{0.15in} \ans{$(1,1)$}
    \hspace{0.15in} \ans{$(3,5)$} \\[3pt]
    (b) $h(-1)=$ \ans{$0$}, \; so $h^{-1}(0)=$ \ans{$-1$}. \hspace{0.25in}
    (c) One marked point is its own reflection. Which? \ans{$(1,1)$} \\[3pt]
    (d) Does $k$ pass the horizontal line test? \ans{no} \; Give a horizontal line that proves
    your answer: \ans{$y=1$ (hits at $x=\pm2$)} \\[3pt]
    (e) Name a restricted domain that would give $k$ an inverse function, and write that inverse.
    \ans{$x\ge0$, giving $k^{-1}(x)=\sqrt{x+3}$ \; (or $x\le0$, giving $-\sqrt{x+3}$)}

  \item \textbf{Error analysis.} Two classmates, two different misunderstandings.
    \begin{center}
    \begin{tcolorbox}[colback=white, colframe=charcoal!45, boxrule=0.7pt, arc=1.5mm,
        width=0.80\linewidth, left=3mm, right=3mm, top=1.5mm, bottom=1.5mm]
    \centering\small
    \renewcommand{\arraystretch}{1.5}
    \begin{tabular}{@{}l l@{}}
    \textbf{Priya writes:} & $f(x)=5x$, \; so \; $f^{-1}(x)=\dfrac{1}{5x}$ \\
    \textbf{Devon writes:} & $f(x)=\dfrac{x}{4}-9$, \; so \; $f^{-1}(x)=4x+9$ \\
    \end{tabular}
    \end{tcolorbox}
    \end{center}
    \vspace{-4pt}
    (a) Give the correct inverse in each case. \; Priya's: \ans{$\frac{x}{5}$} \; Devon's:
    \ans{$4(x+9)=4x+36$} \\[3pt]
    (b) Catch each error with one number. For Priya, $f(3)=$ \ans{$15$}, so $f^{-1}$ must send
    that back to \ans{$3$}; hers sends it to \ans{$\frac{1}{75}$}. \\[3pt]
    \hspace*{0.3in} For Devon, $f(4)=$ \ans{$-8$}, so $f^{-1}$ must send that back to
    \ans{$4$}; his sends it to \ans{$-23$}. \\[3pt]
    (c) Name each mistake in a few words --- they are two \emph{different} misunderstandings. \\[2pt]
    Priya: \ans{read the $-1$ as an exponent and took the reciprocal} \\[3pt]
    Devon: \ans{reversed the operations but not the order --- multiplied before adding}

  \item \textbf{Justify.} Answer both in complete sentences.
    \begin{enumerate}[label=(\alph*), leftmargin=1.9em, itemsep=4pt]
      \item $y=x^{2}$ must have its domain restricted before it can have an inverse function, but
            $y=x^{3}$ does not. Explain the difference, using the horizontal line test. Then say
            what this has to do with the fact that $\sqrt{x}$ accepts only half the number line
            while $\sqrt[3]{x}$ accepts all of it.
      \item Let $f(x)=\sqrt{x+3}$. State the domain and range of $f$, then the domain and range of
            $f^{-1}$ --- and explain how you got the second pair from the first \emph{without}
            finding $f^{-1}$ at all.
    \end{enumerate}
    \ansline{(a) $y=x^{2}$ fails the horizontal line test: $y=4$ crosses it at both $x=-2$ and
    $x=2$, so swapping coordinates would hand the input $4$ two different outputs --- not a
    function.}
    \ansline{\ \ \ \ $y=x^{3}$ passes: it climbs the whole way and never repeats an output, so
    every horizontal line meets it exactly once and the swap is safe.}
    \ansline{\ \ \ \ The domain of an inverse is the range of the original. Restricted to $x\ge0$,
    $y=x^{2}$ has range $y\ge0$, so $\sqrt{x}$ inherits $x\ge0$ as its domain; $y=x^{3}$ has range
    all reals, so $\sqrt[3]{x}$ inherits all reals.}
    \ansline{(b) $f$: domain $x\ge-3$, range $y\ge0$. \; $f^{-1}$: domain $x\ge0$, range
    $y\ge-3$.}
    \ansline{\ \ \ \ An inverse reads every arrow backwards, so the inputs and outputs trade
    places: whatever came \emph{out} of $f$ is exactly what goes \emph{into} $f^{-1}$. Swapping the
    two sets is the entire calculation --- no formula required. (For the record, $f^{-1}(x)=x^{2}-3$
    with $x\ge0$.)}
\end{enumerate}
\end{notesbox}

\begin{extensionbox}
\textbf{Building a clock.} A pendulum of length $L$ feet takes
\[
  T(L)=2\pi\sqrt{\frac{L}{32}} \ \text{ seconds}
\]
to complete one full swing. A clockmaker has the opposite problem: they know the swing time they
\emph{want}, and need the length that produces it.
\begin{enumerate}[label=(\alph*), leftmargin=1.8em, itemsep=5pt]
  \item Complete the table. Round to two decimal places.
    \begin{center}
    \renewcommand{\arraystretch}{1.5}
    \begin{tabular}{@{}|c|c|c|c|@{}}
    \hline
    $L$ (feet) & $2$ & $8$ & $32$ \\
    \hline
    $T(L)$ (seconds) & \ans{$1.57$} & \ans{$3.14$} & \ans{$6.28$} \\
    \hline
    \end{tabular}
    \end{center}
  \item Each length in that table is $4$ times the one before it. What happened to the swing times,
        and which feature of a square root graph made that inevitable?
  \item Build the clockmaker's function by swapping and solving. Show that
        $L(T)=\dfrac{8T^{2}}{\pi^{2}}$.
  \item A grandfather clock is built so that one swing takes exactly $2$ seconds. How long is its
        pendulum, in feet and in inches? (Look up a real grandfather clock afterwards.)
  \item Verify that $T$ and $L$ are inverses by composing them \emph{both} ways.
  \item State the domain and range of $T$ \emph{in context}, then the domain and range of $L$. Why
        does this situation hand you the restriction for free, without you having to decide
        anything?
\end{enumerate}
\ansline{(a) $T(2)=2\pi\sqrt{1/16}=\frac{\pi}{2}\approx1.57$; \; $T(8)=2\pi\sqrt{1/4}=\pi\approx3.14$;
\; $T(32)=2\pi\sqrt{1}=2\pi\approx6.28$.}
\ansline{(b) They \textbf{doubled} each time. Quadrupling the input doubles a square root's output,
because $\sqrt{4L}=2\sqrt{L}$ --- the flattening you first read off the graph in Lesson 6.0.}
\ansline{(c) $T=2\pi\sqrt{L/32} \Rightarrow \frac{T}{2\pi}=\sqrt{\frac{L}{32}} \Rightarrow
\frac{T^{2}}{4\pi^{2}}=\frac{L}{32} \Rightarrow L=\frac{32T^{2}}{4\pi^{2}}=\frac{8T^{2}}{\pi^{2}}$.
\; Squaring is legal here because both sides are non-negative --- no extraneous root (6.5).}
\ansline{(d) $L(2)=\frac{8\cdot4}{\pi^{2}}=\frac{32}{\pi^{2}}\approx3.24$ ft $\approx 38.9$ inches.
A real grandfather clock's pendulum is about $39$ inches --- this is why.}
\ansline{(e) $T\big(L(T)\big)=2\pi\sqrt{\frac{8T^{2}/\pi^{2}}{32}}=2\pi\sqrt{\frac{T^{2}}{4\pi^{2}}}
=2\pi\cdot\frac{T}{2\pi}=T$, and
$L\big(T(L)\big)=\frac{8}{\pi^{2}}\left(2\pi\sqrt{\frac{L}{32}}\right)^{2}
=\frac{8}{\pi^{2}}\cdot4\pi^{2}\cdot\frac{L}{32}=L$.}
\ansline{(f) $T$: domain $L>0$, range $T>0$; $L$: domain $T>0$, range $L>0$ --- the same two sets,
swapped. The context restricts them for free: a pendulum cannot have negative length and a swing
cannot take negative time, so the half-line you had to impose \emph{by hand} on $y=x^{2}$ is
imposed here by the situation.}
\end{extensionbox}

\begin{spiralbox}
\textbf{That is Unit 6 --- and the loop is closed.} Look back at Lesson 6.0. You noticed then that
$y=\sqrt{x}$ gets only half the number line while $y=\sqrt[3]{x}$ gets all of it, that one has an
endpoint and the other does not, that one has an absolute minimum and the other has no extrema at
all. Today you found the single sentence underneath every one of those differences: $x^{2}$ has two
arms and $x^{3}$ has one, so the square must be cut down before it is allowed a partner and the cube
never has to be. Everything in between was the machinery --- rational exponents (6.1) for
\emph{writing} the two families in one notation, simplifying and operating (6.2--6.3), graphing and
transformations (6.4), solving equations and the extraneous roots that squaring creates (6.5), and
composition (6.6), which turned out to be the tool that \emph{proves} an inverse rather than just
suggesting one.
\\[4pt]
\textbf{Next: the Unit 6 test.} Bring back every part of it --- convert between radicals and rational
exponents, simplify and combine radical expressions, rationalize a denominator, read the
characteristics of a radical graph, transform one, solve a radical equation and check for
extraneous solutions, compose two functions and find the domain, and build and justify an inverse.
The practice test is in your packet.
\end{spiralbox}

\begin{teachernote}
\small \textbf{Grading priorities, with the unit test three days away.}
\begin{itemize}[leftmargin=1.3em, nosep]
  \item \textbf{Item 1(c) and item 5 (Devon)} are the same error twice --- right operations, wrong
        order. A student who misses both has not internalized socks-and-shoes and will lose the
        whole inverse strand on the test. Two minutes of Warm-Up item 3 fixes it.
  \item \textbf{Item 2(f) and item 6(b)} are the A2.F.2i justification items. Treat a missing
        restriction as wrong, not incomplete: $(x-4)^{2}$ without $x\ge4$ is a different function
        and it fails the composition test on every input below $4$. The one sentence to demand:
        \emph{the domain of $f^{-1}$ is the range of $f$}.
  \item \textbf{Item 2(c)} is the sleeper. $f(x)=7-2x$ trips students who reach for ``divide by
        $-2$'' first. It is worth putting on the board.
  \item \textbf{Item 6(a) is the capstone question of the unit.} A complete answer connects three
        things students have met on three different days: the horizontal line test, the restricted
        domain, and the domain of $\sqrt{x}$. If most of the class gets two of the three, spend the
        first ten minutes of the review day on that chain rather than on more procedures.
  \end{itemize}
\textbf{The extension} is genuinely worth assigning to everyone, not just the fast finishers: it is
the only place in the unit where the \emph{inverse} is the thing a real person actually needs --- a
clockmaker never asks how fast a given pendulum swings, they ask how long to make it. Part (d)
producing $\approx39$ inches, which students can then check against any real grandfather clock, is
the payoff. Part (f) is the quiet one: every restricted domain in this unit had to be argued for,
and here the situation supplies it without anyone deciding anything --- the same move as Lesson
6.0's free-fall and skid-mark models.
\end{teachernote}

\end{document}
